\begin{problem}[The principal-open localization theorem]
Let $f\in A$.  Prove that extension and contraction identify
\[
\Spec A_f
\]
homeomorphically with
\[
D_A(f)\subseteq\Spec A.
\]
\end{problem}

\paragraph{Why this problem matters.}
This theorem is the algebraic heart of principal opens: restricting geometry to $D(f)$ is
exactly the same as inverting $f$.

\paragraph{Useful ideas before starting.}
Use the prime correspondence for localization at
$S=\{1,f,f^2,\dots\}$ and then compare basic opens on both sides.

\paragraph{Guided hints.}
A prime $\mathfrak p$ is disjoint from $S$ iff $f\notin\mathfrak p$.  Under the resulting
bijection, determine the image of $D_{A_f}(g/1)$.

\begin{solution}
The localization prime correspondence gives a bijection between primes of $A_f$ and primes
$\mathfrak p$ of $A$ disjoint from $S$.  Since $S$ consists of powers of $f$, disjointness
is equivalent to $f\notin\mathfrak p$.  Thus extension and contraction give inverse
bijections
\[
\Spec A_f\longleftrightarrow D_A(f).
\]

To compare topologies, let $g\in A$.  A point $\mathfrak pA_f$ belongs to
$D_{A_f}(g/1)$ iff $g/1\notin\mathfrak pA_f$.  Under contraction this is equivalent to
$g\notin\mathfrak p$.  Together with $f\notin\mathfrak p$, the corresponding set in
$\Spec A$ is
\[
D_A(f)\cap D_A(g)=D_A(fg).
\]
These sets form bases for the source and the subspace topology on $D_A(f)$.  Hence the
bijection is a homeomorphism.
\end{solution}

\paragraph{What happened mathematically.}
Localization and principal-open restriction are the same operation viewed from opposite
sides of the algebra/geometry dictionary.

\paragraph{Extensions.}
Show that the universal property of $A_f$ is exactly the universal property of forcing $f$
to be invertible.

\paragraph{Provenance.}
Canonical expansion of the localization theorem in the legacy VI/08 section of
\emph{Theory of Algebraic Geometry 1}.
